% ========================================================================================
\chapter{About}\label{ch:One}
This document demonstrates the appearance of a thesis build with the \texttt{\classname}
class. A full description will be available soon at my
\href{https://timroith.github.io/}{webpage}.
If you have any questions or found a problem/bug please don't hesitate to contact me
via e-mail\footnote{tim.roith@fau.de}.
% ========================================================================================
\section{Tiny Chickens}\label{sec:tiny-chickens}

\begin{figure}[h]
    \centering
    \includegraphics[width=0.5\textwidth]{figures/tiny_chicken}
    \caption{A tiny chicken.}
    \label{fig:tiny-chicken}
\end{figure}

Tiny chickens, commonly known as \emph{bantam} chickens, are miniature versions of
standard chicken breeds. They are popular for backyard keeping due to their small size,
docile temperament, and ornamental appearance. The following is an overview of the
major types of tiny chickens:

\begin{enumerate}[label=(\roman*)]
    \item \textbf{True Bantams} --- breeds with no large fowl counterpart, naturally
          small by origin. Examples include the Sebright, Rosecomb, and Nankin.
    \item \textbf{Miniaturized Bantams} --- scaled-down versions of standard breeds,
          selectively bred to be roughly one-quarter to one-fifth the size of the
          original. Examples include the Miniature Rhode Island Red and Miniature
          Leghorn.
    \item \textbf{Developed Bantams} --- breeds developed over time from crosses of
          true bantams and miniaturized bantams. The Pekin (Cochin Bantam) is a
          prominent example.
    \item \textbf{Frizzle Bantams} --- characterized by feathers that curl outward
          rather than lying flat, giving a distinctive fluffy appearance. The Frizzle
          can appear in many bantam breeds.
    \item \textbf{Silkie Bantams} --- known for their unusually soft, silk-like plumage
          and black skin. They are among the most popular ornamental bantam breeds
          worldwide.
    \item \textbf{Belgian Bantams} --- a group of true bantams originating from Belgium,
          including the Barbu d'Anvers (Antwerp Belgian) and Barbu de Watermael,
          notable for their bearded and muffed appearance.
\end{enumerate}
